\section{Related Literature}\label{sec:lit}

This paper sits at the intersection of two literatures.\footnote{\textcolor{red}{\textbf{This paper is a work in progress intended solely for feedback and should not be circulated.}}} The first concerns gender, voice, and communication. \cite{gilligan1993different} provides the motivating intuition that dominant evaluative frameworks may overlook alternative ways of presenting judgment and reasoning. Although that literature is not about economics articles specifically, it suggests that apparently neutral standards of professional communication may embed preferences for particular styles of exposition.

The second is the economics literature on gender and the production and evaluation of research. \cite{card2020} study whether referees and editors in economics treat submissions differently by gender, and \cite{sarsons2017recognition} shows that collaborative work can be differentially recognized by gender. Most directly, \cite{hengel2022publishing} documents that female-authored economics papers are more readable and that the writing gap widens during peer review, consistent with women facing higher writing standards in publication.

This paper contributes to that literature in three ways. First, rather than studying editorial decisions, citations, or promotion outcomes, we focus on writing itself as an observable input into evaluation. Second, we move beyond a single readability measure by using a multidimensional rubric that distinguishes clarity and structure from tone, assertiveness, and affect. This matters because several common narratives about gendered writing, such as those centered on hedging, confidence, and emotionality, receive limited support in our data. Third, by organizing the evidence in three steps, validation of the LLM measure, unconditional mean differences, and controlled regressions, we show both that the measurement is anchored to existing work and that the main patterns are not driven solely by compositional differences across journals, years, or paper types.